% titel: Blatt 1, A4: Protonen, Neutronen, Elektronen
% Teil II: der durchgerechnete Weg.

\subsection*{Gegeben und gesucht}
Protonen ($Z$), Neutronen ($N=A-Z$), Elektronen ($Z-q$) für fünf Teilchen.

\subsection*{Rechnung}
\begin{center}
\begin{tabular}{llll}
$\mathrm{{}^{16}O}$    & $Z{=}8,\  N{=}8,\  e^{-}{=}8$    & & $N/Z=1.00$\\
$\mathrm{{}^{31}P}$    & $Z{=}15,\ N{=}16,\ e^{-}{=}15$ & & $N/Z=1.07$\\
$\mathrm{{}^{63}Cu^{2+}}$ & $Z{=}29,\ N{=}34,\ e^{-}{=}27$ & & $N/Z=1.17$\\
$\mathrm{{}^{107}Ag^{+}}$ & $Z{=}47,\ N{=}60,\ e^{-}{=}46$ & & $N/Z=1.28$\\
$\mathrm{{}^{238}U}$   & $Z{=}92,\ N{=}146,\ e^{-}{=}92$ & & $N/Z=1.59$
\end{tabular}
\end{center}

\begin{ergebnis}
Tabellarisch wie oben; Kernzusammensetzung durch Ionenladung unverändert,
nur die Elektronenzahl ändert sich ($\mathrm{Cu^{2+}}$: $29-2=27$).
\end{ergebnis}

\subsection*{f) Protonen/Neutronen-Verhältnis}
Bei leichten Kernen ist $N/Z\approx1$; mit zunehmender Ordnungszahl steigt
$N/Z$ (U: $1.59$). Grund: mit mehr Protonen wächst die elektrostatische
Abstossung, die durch zusätzliche Neutronen (nur starke Kraft, keine Ladung)
kompensiert wird.

\begin{kontrolle}
Für stabile Kerne gilt umso mehr Neutronen, desto schwerer das Element ---
bekanntes Bild der Stabilitätslinie. Die Rechnung ist leer von Einheiten
(Anzahl).
\end{kontrolle}

\begin{merksatz}
Aus dem Nuklidsymbol: $Z$ festlegen, $N=A-Z$, Elektronen $=Z-q$. Kationen
verlieren Elektronen, nie Nukleonen.
\end{merksatz}
